%% ma: L12-B01-0011
%% lop: 12
%% chuong: 1
%% bai: 1
%% dang: 12.1.7
%% loai: TN4
%% muc: VD
%% dap_an: "B"
%% nguon: "(NBV)-ĐỀ SỐ 7-MỖI NGÀY 1 ĐỀ THI 2025.pdf (D07, MỖI NGÀY 1 ĐỀ - Đề số 7), Phần I Câu 12"
%% kien_thuc: "Bài 1 (điều kiện hàm số nghịch biến trên R), dấu tam thức bậc hai (Lớp 10 Bài 17)"
%% trang_thai: cho-duyet
%% ly_do: "Dạng toán mới đề xuất 12.1.7 (Số giá trị nguyên của tham số m để hàm bậc ba nghịch biến trên R); chờ giáo viên duyệt dạng."
%% ghi_chu: "Đáp án gốc B đúng (kiểm bằng sympy)."
\begin{ex}
Có bao nhiêu số nguyên $m$ để hàm số $y=-x^3-3(m+1)x^2+3(m+1)x-1$ nghịch biến trên $\mathbb R$?
\choice{$3$}{\True $2$}{$1$}{$0$}
\loigiai{$y'=-3x^2-6(m+1)x+3(m+1)$. Hàm số nghịch biến trên $\mathbb R$ khi và chỉ khi $y'\le0$ với mọi $x$. Vì hệ số của $x^2$ là $-3<0$, điều kiện là $\Delta'=9(m+1)^2+9(m+1)\le0\Leftrightarrow(m+1)(m+2)\le0\Leftrightarrow-2\le m\le-1$. Có hai số nguyên $m=-2$ và $m=-1$.}
\end{ex}
